% origin.tex — "Origin: Playing a Robot Orchestra as One Instrument",
% \input from septagraph.tex. ~half a page. Honest narrative: grid for
% robots -> what did not work -> the protocol. Draft for emgor to re-voice.
% ACTION(emgor): confirm authorship/acknowledgment with Kapur and Colton
%   before submission (the ensemble, its server lineage and its current
%   maintainer are described here; no individuals are named as authors).

Septagraph did not start as a protocol. It started as a grid.

The neuralGrid was designed in mid-2026 as a Monome-style controller
with a second job: a small local model would watch how the player
played and rewrite the grid's own layout, every change quantized to a
shared clock so the model's latency would land on a downbeat and never
be felt. Its first target was the Machine Lab. The ensemble is driven
from a Mac mini running a ChucK server whose lineage runs through
several generations of the lab's students and faculty; composers send
OSC to one port, one address per robot, and the server turns each
message into MIDI or serial commands for that machine. The obvious
first build was a direct script: grid cells hand-mapped to robots, a
pad press becoming an OSC message to the lab server. That script still
exists and is still the gig tool; it is also the end-to-end proof that
the hardware works.

Several things about the first design did not survive contact with the
room. The four ribbon sensors on the grid's edges were originally
menus --- touch one to select a mode, then scroll it --- and the
player's last-touched ribbon decided what the grid meant. In practice a
menu on a performance surface meant looking at the instrument instead
of playing it, and ``touched last'' turned sensor noise into mode
changes. The ribbons were demoted to pure modulators, one axis each,
and menus moved to a web page. Beat-quantized reconfiguration went the
same way: the robots share no clock, and the compositions written for
them assume none, so reconfiguring on a beat meant importing a
metronome nobody in the room was listening to. Reconfiguration now
lands at detected gesture boundaries --- moments of stillness in the
input stream --- which need no clock and are, in practice, the moments
when a change of mapping is least disruptive.

The robots also forced the shape of the renderer manifest. A
synthesizer can be told to reach any point in a continuous region; a
robot cannot. Each beater on each machine is one fixed sound at one
fixed place in the room, so the ensemble's reachable set is a finite
catalog of measured positions, not a bound. A catalog of positions is
still not enough: a composer who sends the same beater too many hits
per second gets softer strikes and, on GanaPati, a thermal shutdown.
So every catalog point carries its mechanical constraints --- maximum
re-strike rate and a force range --- and a controller that respects
them will never ask a robot for a roll it cannot play. This is the
part of the protocol we would not have designed for synthesizers
alone, and it is what makes a physical instrument a first-class
renderer rather than a special case.

Once the grid's output had become ``a position in a perceptual space,
with a tolerance'' rather than ``robot $n$, beater $m$'', the
separation described in the introduction was already in place: the
same stream that drives the robots drives a browser synthesizer, and
the instrument on the table no longer knows or cares which one is
listening. The rest of this paper describes that stream.
